\section{Motivation}

Tool-using LLM agents operate under a persistent tension between \emph{answer quality} and \emph{execution cost}. In principle, additional reasoning, retrieval, and verification steps may improve the final answer. In practice, however, these additional steps are not free: they increase token consumption, wall-clock latency, and trajectory length, and they also enlarge the accumulated prompt context during iterative execution. As a result, practical agent design is not only a reasoning problem, but also a control problem.

This tension is especially visible when contrasting two dominant design patterns. Fixed workflows provide bounded and predictable execution cost, but they are often too rigid to adapt to the difficulty of the current instance or to changing intermediate evidence. On the other hand, free-form multi-step agents, such as ReAct-style systems, can adapt their behavior online, but they may also over-execute. Typical failure modes include repeated retrieval after useful evidence has already been obtained, redundant tool calls with highly similar intent, and continued reasoning even when the marginal value of another step is low. These behaviors do not merely waste computation; they also make agent trajectories harder to interpret and harder to control.

This observation motivates the central perspective of our work: \emph{agent orchestration should be treated as an explicit decision problem rather than an implicit side effect of prompting}. Instead of asking only whether a model can reason or call tools, we ask a more operational question: \emph{when is another action actually worth taking?} Under this view, a tool-using agent should not only know how to retrieve or verify, but should also decide whether retrieval or verification is justified at the current step.

Our goal is therefore different from simply maximizing final answer quality. We aim to make the \emph{quality--cost trade-off} explicit, controllable, and analyzable. This naturally leads to a utility-guided formulation in which candidate actions are evaluated not only by their potential contribution to answer quality, but also by the expected cost of continuing, the current level of uncertainty, and the risk of redundant behavior. Such a formulation turns orchestration into a lightweight control layer between the language model and the external environment.

This perspective is useful for at least three reasons. First, it provides a more interpretable view of multi-step agent execution: the continuation of the trajectory is no longer governed solely by prompt behavior, but also by explicit control signals. Second, it is relevant in budget-sensitive deployments, where excessive reasoning steps directly translate into monetary and latency costs. Third, it matters in context-limited settings, because unnecessary intermediate steps accelerate prompt growth and increase the burden on the context window even when hard truncation is not triggered.

Taken together, these considerations suggest that efficient LLM agents require more than strong reasoning ability or access to external tools. They also require a decision mechanism that regulates \emph{when to continue}, \emph{what action to take next}, and \emph{when to stop}. This is the role played by the orchestration policy studied in this paper.